% ID: SEL-D0060-OB-E04
% ESTADO: cotejada. Segunda pasada visual y borrador compilado revisados.
% Fuente: Selectividad/sel_2011_matematicas_aplicadas/223A-EXAMEN 3.pdf
% SHA-256: 5998699e0b2992b988a8e0ed0d2e3581dd8fac52d6f9802587ebd40fee5ef608
% Página: 2; opción B; ejercicio original 4.

Con el fin de estudiar el peso medio de los perros recién nacidos de una determinada raza, se tomó una muestra en una clínica veterinaria y se obtuvieron los siguientes pesos, medidos en kg:
\[
1.2\quad0.9\quad1\quad1.2\quad1.1\quad1\quad0.8\quad1.1
\]
Se sabe que el peso de los cachorros de esta raza se distribuye según una ley Normal con desviación típica $0.25$ kg.
\begin{enumerate}[label=\alph*)]
\item \textbf{(1.5 puntos)} Obtenga un intervalo de confianza para estimar la media poblacional, al $95\%$.
\item \textbf{(0.5 puntos)} Halle el error máximo que se cometería usando el intervalo anterior.
\item \textbf{(0.5 puntos)} Razone cómo variaría la amplitud del intervalo de confianza si, manteniendo el mismo nivel de confianza, aumentásemos el tamaño de la muestra.
\end{enumerate}
